% ---
% titre: Problème de proportionnalité- les séances de cinéma
% theme: FA
% chapitre: Proportionnalité
% sous_chapitre: Résoudre des problèmes
% niveau: [10e]
% type: EVAL
% tags: [c-proportionnalite,p-question-TS]
% difficulte: 1
% corrige: true
% ---

\question
Au cinéma Pathé, une séance coûte en moyenne 17.-.
\begin{parts}
\vspace{10pt}\part[1] Théo décide d'aller au cinéma une fois par semaine pendant les 38 semaines d'école de cette année. Quel budget total doit-il prévoir ?

\vspace{10pt}{\footnotesize\raggedleft Espace pour ta démarche et tes calculs\par}
\begin{solutionorgrid}[70pt]

\vspace{5pt} Budget total: 38 \cdot 17 = 646,- \hfill\textbf{(0.5pt)}
\end{solutionorgrid}

\begin{tcolorbox}[enhanced jigsaw, interior hidden, colframe=box, parbox=false]%
Réponse: 
\sol{\vspace{20pt}}{Il doit prévoir env 650.- de budget total. \textbf{(0.5pt)}}
\end{tcolorbox}


\vspace{20pt}\part[1] Son amie Jenna aimerait l'accompagner, mais ses parents ne lui accordent que 150.- pour l'année. Combien de séances pourra-t-elle s'offrir ?

\vspace{10pt}{\footnotesize\raggedleft Espace pour ta démarche et tes calculs\par}
\begin{solutionorgrid}[70pt]

\vspace{5pt}Nombre de séances: 150 : 17 = 8,8 \hfill\textbf{(0.5pt)}
\end{solutionorgrid}

\begin{tcolorbox}[enhanced jigsaw, interior hidden, colframe=box, parbox=false]%
Réponse: 
\sol{\vspace{20pt}}{Elle pourra s'offrir 8 séances. \textbf{(0.5pt)}}
\end{tcolorbox}
\end{parts}

